\section*{Annex C --- Debrief Form}
\addcontentsline{toc}{section}{Annex C --- Debrief Form}

The following questionnaire was administered to participants after the first playtest. Questions aim to gather data on entertainment value, format distinctiveness, participant intentions, and transformative effects.

\begin{enumerate}
    \item How fun was the debate? (scale 1--10)
    \item How fun was it compared to a debate in a traditional format? (5-choice scale)
    \item Was there anything surprisingly entertaining or engaging about it?
    \item Was there anything about the format that was quite difficult, to the point of frustration?
    \item Was there anything about the rules of the format that was confusing?
    \item Would you play this format some more? (Yes / No)
    \item If you would play it some more, in what circumstance would you?
    \item If there was to exist an alternative debate circuit centred around this format, would you attend those competitions? Even at the expense of competitions in the WS or BP format?
    \item Do you believe you played differently than you would have in a debate in WS/BP with a rather similar motion? If so, how?
    \item Do you feel different than how you usually feel after a debate? If so, how?
    \item Did you perceive yourself and others speaking similarly or differently than how you do in a regular debate? If so, how?
    \item Were you to redo this debate, what would you do differently?
    \item Were you to play more in this format, but on different motions, what would you do differently?
    \item What do you believe you learned from this debate?
    \item What do you believe this debate format does better, as a learning/empowerment tool, than WS or BP? Do you believe it does something worse?
\end{enumerate}

%%%%%%%%%%%%%%%%%%%%%%%%%%%%%%%%%%%%%%%%%

\newpage
\section*{Annex D --- Playtest Motion}
\addcontentsline{toc}{section}{Annex D --- Playtest Motion}

The following motion was used in the first playtest at the Jean Monet debate club, Bucharest, March 2026.

\paragraph{First version.} Used in the first playtest at the Jean Monet debate club, Bucharest, March 2026.

\begin{quote}
You are Alex. You are a lawyer. You are close to your older brother Matei. He is a very caring person to everyone, always eager to help even strangers. He always helped you when you were in need. However he was always quite dishonest with you, and you caught him lying even about important things. Matei is also an alcoholic and you believe he is quite bigoted.

\medskip

Matei was involved in a hit and run recently. He hit a cyclist that ended up hospitalised. The cyclist remembered his license plate only partially and he is now being investigated by the police. From your knowledge they only have circumstantial evidence, but he still risks going to jail. He asked you to lie that he was with you at home at the time of the accident, in order to provide him an alibi. Your assessment as a lawyer is that this will almost certainly exonerate him. However, a very small risk remains that additional material evidence could be found, at which point you might become an accessory to the crime.

\medskip

\textbf{What should you do?}

\medskip

Three Values that you hold, relevant for this decision, are:
\begin{itemize}
    \item You believe that helping people in need is very important, and a moral duty.
    \item You believe in justice: that people should receive what they deserve.
    \item You believe that family should stick together and have each other's back.
\end{itemize}
\end{quote}

\paragraph{Second version.} Used in the second playtest, April 2026.

\begin{quote}
Alex is a criminal defence lawyer with seven years of experience. The eldest of three siblings from a small town, Alex was the first in the family to attend law school. After graduating, Alex spent two years at a legal aid clinic representing clients who could not afford private counsel, before joining a private firm. Alex's first major case there was defending a driver charged with causing serious injury in a road accident --- Alex secured his acquittal.

\medskip

Matei is Alex's older brother and the two have remained close throughout their lives. Matei is known among everyone who knows him for being eager to help others, often going out of his way even for strangers. He helped Alex through several difficult periods. He has also been dishonest with Alex repeatedly over the years, including about important matters. Matei struggles with alcohol and holds a number of views that Alex finds bigoted.

\medskip

Matei was recently involved in a hit-and-run: he struck a cyclist who was hospitalised. The cyclist recalled his license plate only partially, and Matei is now under police investigation. As far as Alex knows, the evidence against him is circumstantial. Matei has asked Alex to lie and say he was at home with Alex at the time of the accident. Alex's legal assessment is that this would almost certainly exonerate him, though a small risk remains that additional material evidence could surface --- at which point Alex would become an accessory to the crime. \textbf{What should you do?}

\medskip

Values:
\begin{enumerate}
    \item You believe that helping people in need is very important, and a moral duty.
    \item You believe in justice, that people should receive what they deserve.
    \item You believe that family should stick together and have each other's back.
\end{enumerate}
\end{quote}

%%%%%%%%%%%%%%%%%%%%%%%%%%%%%%%%%%%%%%%%%

\newpage
\section*{Annex E --- Informed Consent Form}
\addcontentsline{toc}{section}{Annex E --- Informed Consent Form}

\bigskip

I, \underline{\hspace{7cm}} of \underline{\hspace{1.5cm}} years of age, voluntarily accept to participate in a playtest followed by an interview session, both moderated by Andrei Dan Olaru, as part of his \textit{``Applying Transformative Game Design Theory to Competitive Debate''} study. This study seeks to explore the creation of an alternative debate format, informed by modern transformative game design theory. I understand that:

\begin{enumerate}
    \item This session implies video recording of the entire debate playtest and the subsequent collective interview, plus some additional questions using an online questionnaire.
    \item Pseudonyms will be used when analysing the data for confidentiality, and the researchers commit to not revealing the identities of the participants. Individual raw data may be discussed in class or with the supervisor without the participant's name to trace back to. The data obtained will be analysed by grouping and some textual examples will be used anonymously.
    \item The recording will be safely stored in the university's OneDrive to maintain the saved raw interviews protected. The video will be erased one month after the dissertation's results. The transcripts will be erased up to one year after the dissertation's results.
    \item None of the questions will go against my integrity, no harm will come to me during this session, no answer is right or wrong and I may choose not to answer any question(s) without consequences of any kind.
    \item I am free to walk out of the playtest and interview at any time if I deem it necessary without prejudice or any consequences.
\end{enumerate}

For any additional information or difficulty, the participant may contact the researcher Andrei Dan Olaru at \texttt{andrei-dan.olaru.0744@student.uu.se}.

\bigskip\bigskip

\noindent
\begin{tabular}{p{6cm} p{1cm} p{6cm}}
\underline{\hspace{6cm}} & & \underline{\hspace{6cm}} \\
Participant's signature & & Researcher's signature \\
\end{tabular}

\bigskip\bigskip

I, \underline{\hspace{7cm}} have been informed of the conditions and freely accept to participate in the research by Andrei Dan Olaru.

\bigskip

\noindent
\begin{tabular}{p{6cm} p{5cm}}
\underline{\hspace{6cm}} & \underline{\hspace{3cm}} of March 2026 \\
Participant's signature & Date \\
\end{tabular}

%%%%%%%%%%%%%%%%%%%%%%%%%%%%%%%%%%%%%%%%%

\newpage
